\documentclass[11pt]{article}
\usepackage[margin=1in]{geometry}
\usepackage{amsmath,amssymb,booktabs,hyperref}
\usepackage[T1]{fontenc}
\newcommand{\Oh}{\mathrm{Oh}}
\title{Faithful two-pole closure for the viscous drop modes\\
{\large Reid's arbitrary-$\Oh$ eigenvalues mapped onto the drop-mode ODE}}
\author{SpectralKM derivations}
\date{}

\begin{document}
\maketitle

\section{What is being matched, and why}

Each droplet surface mode in this model is carried by a second-order ODE
(\texttt{eq:lambda-l}; implemented in \texttt{src/affine.jl:drop\_affine})
\begin{equation}
  \ddot\beta_l + 2\lambda_l\dot\beta_l + \omega_l^2\beta_l = F_l\,b_l\text{,}
  \label{eq:mode-ode}
\end{equation}
whose characteristic polynomial is $s^2 + 2\lambda_l s + \omega_l^2$. The published choice
takes Lamb's (1881) small-viscosity asymptotics,
\begin{equation}
  \lambda_l^{\text{Lamb}} = \Oh\,(l-1)(2l+1)\text{,}\qquad
  \big(\omega_l^{\text{Lamb}}\big)^2 = \omega_{l,0}^2 = l(l-1)(l+2)\text{,}
  \label{eq:lamb}
\end{equation}
with $\omega_{l,0}$ the inviscid Rayleigh--Lamb capillary frequency. Reid (1960) instead
gives the \emph{exact} linear viscous eigenvalues at arbitrary $\Oh$, as roots
$\sigma$ of a transcendental characteristic equation (see \texttt{src/reid.jl}); with
$\sigma$ a decay rate, perturbations behave as $e^{-\sigma t}$, so the ODE roots are
$s=-\sigma$.

This note fixes the remaining question: \emph{given a pair of Reid eigenvalues, what are the
faithful $(\lambda_l,\omega_l^2)$?} The algebra is verified symbolically in
\texttt{derivations/cas\_reid\_two\_pole.py} (all assertions pass).

\section{The mapping is Vieta, and needs no case split}

For a monic quadratic with roots $s_1,s_2$,
\begin{equation}
  (s-s_1)(s-s_2) = s^2 + 2\lambda_l s + \omega_l^2
  \quad\Longleftrightarrow\quad
  \boxed{\ \lambda_l = -\tfrac{1}{2}(s_1+s_2)\text{,}\qquad \omega_l^2 = s_1 s_2\ }
  \label{eq:vieta}
\end{equation}
exactly, for \emph{any} pair. Two consequences.

\paragraph{Underdamped ($s = -\gamma \pm i\omega$).}
\begin{equation}
  \lambda_l = \gamma\text{,}\qquad \omega_l^2 = \gamma^2 + \omega^2\text{,}
  \label{eq:case-a}
\end{equation}
\emph{not} $\omega^2$. The distinction matters because \eqref{eq:mode-ode} oscillates at
$\sqrt{\omega_l^2-\lambda_l^2}$, and only \eqref{eq:case-a} returns $\omega$. Lamb
\eqref{eq:lamb} carries the same $O(\lambda^2/\omega^2)$ slip, but there it sits inside its
own asymptotic error; here it would be gratuitous.

\paragraph{Overdamped ($s=-\gamma_1,-\gamma_2$, real, $\gamma_1<\gamma_2$).}
\begin{equation}
  \lambda_l = \tfrac{1}{2}(\gamma_1+\gamma_2)\text{,}\qquad
  \omega_l^2 = \gamma_1\gamma_2\text{,}\qquad
  \lambda_l^2-\omega_l^2 = \left(\tfrac{\gamma_1-\gamma_2}{2}\right)^2 \ge 0\text{,}
  \label{eq:case-b}
\end{equation}
so the ODE is genuinely overdamped, degenerating only when the rates coincide. As
$\gamma_2\to\gamma_1$ this reduces continuously to the critically damped case, hence
\eqref{eq:vieta} is a strict generalisation.

\section{The superseded choice, and why it was badly wrong}

An earlier revision of \texttt{drop\_viscous\_coeffs} set, for a real root,
\begin{equation}
  \lambda_l = \gamma_{\text{tracked}}\text{,}\qquad \omega_l^2 = \lambda_l^2\text{,}
  \label{eq:superseded}
\end{equation}
a double root at the tracked rate. This discards $\gamma_2$ and invents a second rate. Two
compounding errors made it worse than a mild approximation:

\begin{enumerate}
\item \textbf{Wrong branch.} Continuation in $\Oh$ from the complex pair does \emph{not} land
on the slowest real root once the pair has merged. Measured at $l=2$ by scanning the real
axis (the residual is real for real $s>0$, since $q=\sqrt{s/\Oh}$ is then real and so is
$j_l$):

\begin{center}
\begin{tabular}{llll}
\toprule
$\Oh$ & real roots on $(10^{-4},200)$ & tracked & slowest \\
\midrule
$0.8$ & $2.040,\ 3.670,\ 23.94,\ 63.04,\ 118.3,\ 189.4$ & $3.670$ & $2.040$ \\
$1.0$ & $1.280,\ 5.846,\ 29.93,\ 78.80,\ 147.9$ & $5.846$ & $1.280$ \\
$3.0$ & $0.357,\ 20.97,\ 89.86$ & $20.97$ & $0.357$ \\
\bottomrule
\end{tabular}
\end{center}

Continuation consistently returns the \emph{second} root. The slowest root decreases with
$\Oh$ (as a very viscous drop must relax ever more slowly, rate $\propto 1/\Oh$), whereas the
tracked one grows --- the diagnostic that exposed this.

\item \textbf{Wrong stiffness.} Applying \eqref{eq:vieta} to the two slowest roots gives
$\omega_l^2 = 7.487,\ 7.485,\ 7.480$ at $\Oh=0.8,1.0,3.0$, against
$\omega_{2,0}^2 = l(l-1)(l+2) = 8$. The faithful mapping therefore \textbf{preserves the
capillary stiffness}, essentially independent of $\Oh$, as it must: viscosity damps, it does
not change what restores. The superseded choice \eqref{eq:superseded} instead gave
$\omega_l^2=\lambda_l^2 = 20.97^2 \approx 440$ at $\Oh=3$, a stiffness too large by a factor
$59$.
\end{enumerate}

That $\omega_l^2\approx\omega_{l,0}^2$ is not an input to the construction; it emerges, and so
serves as a cheap correctness check on the root pair.

\section{Where the overdamped branch is actually reached}

The capillary pair merges onto the real axis (mode becomes aperiodic) at an $\Oh$ that falls
with $l$: measured between $0.5$ and $0.8$ for $l=2,4$, and between $0.3$ and $0.5$ for
$l=8,16$. So the overdamped case is not hypothetical for $\Oh\gtrsim0.5$, and high modes enter
it first.

\section{Algorithm}

For each $l$, at the requested $\Oh$:
\begin{enumerate}
  \item Track the least-damped eigenvalue by continuation in $\Oh$ from a value small enough
        that Lamb is genuinely asymptotic (\texttt{reid\_root\_tracked}).
  \item If $|\operatorname{Im}\sigma|>0$: the pair is $\sigma,\bar\sigma$; apply
        \eqref{eq:case-a}.
  \item If $\operatorname{Im}\sigma=0$: the pair has merged. Scan the positive real axis for
        roots (real residual, bisect on sign changes, rejecting pole crossings at zeros of
        $j_l$), take the \emph{two slowest}, apply \eqref{eq:case-b}. Do \emph{not} rely on
        the continued root being the slowest.
  \item Sanity-check $\omega_l^2$ against $\omega_{l,0}^2$; a large discrepancy indicates a
        bad root pair rather than physics.
\end{enumerate}

\section{Scope}

Two limitations are inherited and unchanged by this note. First, \eqref{eq:mode-ode} is a
two-pole model: Reid's spectrum also contains an infinite diffusive ladder (internal Stokes
relaxations, rates $\approx\Oh\,q_j^2$ with $q_j$ the zeros of $j_l$), so the \emph{forced}
response carries memory that no second-order ODE reproduces. Matching two poles exactly is
right for free decay and approximate under forcing; the $n$-pole generalisation is deferred.
Second, the residues are not used here --- only the poles --- so this note says nothing about
the amplitude with which pressure forcing excites each mode.

\end{document}
